\section{A proper central Frattini lifting step}

We prove the cyclic-kernel case of the construction in
\cite{Witt1936}, pp.~238--240, in the form used by
\cite{Camina1997}, pp.~103--104. Witt treats a central elementary
abelian kernel; Camina's remark (ii) permits successive kernels of
order $p$. We use Camina's Frattini hypothesis and express the
properness argument through maximal subgroups. The two additive
equations are supplied by the explicit constructions above.

\begin{lemma}[A Frattini subgroup has no proper supplement]
\label{lem:frattini-no-proper-supplement}
\notready
Let $E$ be a finite group. Write $\Phi(E)$ for the intersection
of its maximal proper subgroups. If $N\le\Phi(E)$ and $H\le E$
satisfy $NH=E$, then $H=E$.
\end{lemma}

\begin{proof}
Suppose $H$ is proper. The set of proper subgroups containing
$H$ is nonempty and finite. A maximal member $M$ of this set
is a maximal proper subgroup of $E$: any larger proper subgroup
would still contain $H$ and contradict its choice.
Then $N\le\Phi(E)\le M$ and $H\le M$, so $NH\subseteq M$.
This contradicts $NH=E$ and $M\ne E$.
\end{proof}

Fix a prime $p$, a field $F$ of characteristic $p$, and a finite
Galois extension $K/F$ with $G=\operatorname{Gal}(K/F)$.
For the central-extension statements below, let $E$ be a finite group and
$q:E\to G$ a surjective homomorphism. Assume that
$N=\ker q$ is central, has order $p$, and lies in $\Phi(E)$.
Choose an isomorphism $\iota:(\mathbf F_p,+)\to N$, with
its values also regarded as elements of $E$, and a section
$u:G\to E$ of $q$ with $u_1=1$.
The inclusion of $\mathbf F_p$ into $K$ is understood whenever
prime-field values occur in a field equation. Group actions on
fields are on the left, so $(st)(x)=s(t(x))$.

\begin{lemma}[The factor set of a central section]
\label{lem:central-extension-factor-set}
\source{witt1936:page-0002,camina1997:page-0003}
\lean{Camina1997.Witt.factorSet_spec, Camina1997.Witt.factorSet_one_left, Camina1997.Witt.factorSet_one_right, Camina1997.Witt.factorSet_cocycle, Camina1997.Witt.centralCoordinate_unique, Camina1997.Witt.centralCoordinates, Camina1997.Witt.central_coordinates_mul}
\leanok
There is a unique function $c:G\times G\to\mathbf F_p$ with
\[
 u_su_t=\iota(c(s,t))u_{st}.
\]
It is a normalized two-cocycle:
$c(1,s)=c(s,1)=0$ and
\[
 c(s,t)+c(st,r)=c(t,r)+c(s,tr).
\]
Every element of $E$ has a unique expression $\iota(a)u_s$,
and multiplication in these coordinates is
\[
 (\iota(a)u_s)(\iota(b)u_t)
   =\iota(a+b+c(s,t))u_{st}.
\]
\end{lemma}

\begin{proof}
The element $u_su_tu_{st}^{-1}$ belongs to $\ker q=N$.
Existence and uniqueness of $c(s,t)$ follow from bijectivity
of $\iota$ onto $N$. Setting either argument equal to $1$
and using $u_1=1$ gives normalization.
Compute $(u_su_t)u_r$ and $u_s(u_tu_r)$ using the defining
formula. Since $N$ is central, their respective expressions are
\[
 \iota(c(s,t)+c(st,r))u_{str},\qquad
 \iota(c(t,r)+c(s,tr))u_{str}.
\]
Canceling $u_{str}$ and using injectivity of $\iota$ gives
the two-cocycle identity.
For $e\in E$, put $s=q(e)$. Then $eu_s^{-1}\in N$, giving
$e=\iota(a)u_s$ for a unique $a$. Applying $q$ to any such
expression recovers $s$, so both coordinates are unique.
Finally centrality permits moving $\iota(b)$ past $u_s$,
and the section formula gives the asserted multiplication.
\end{proof}

\begin{lemma}[The Frattini hypothesis makes the parameter nontrivial]
\label{lem:witt-frattini-parameter-not-image}
\uses{lem:frattini-no-proper-supplement,lem:central-extension-factor-set}
\source{witt1936:page-0004,camina1997:page-0004}
\notready
Suppose $\delta_s\in K$ and $\gamma\in K$ satisfy
\[
 \delta_1=0,\qquad
 \delta_s+s(\delta_t)-\delta_{st}=c(s,t),\qquad
 s(\gamma)-\gamma=\delta_s^p-\delta_s.
\]
Then $\gamma\notin\wp(K)$, where $\wp(x)=x^p-x$.
\end{lemma}

\begin{proof}
Suppose $\gamma=\wp(\beta)$ with $\beta\in K$.
Put $a_s=\delta_s-(s(\beta)-\beta)$. Additivity of $\wp$
and its compatibility with automorphisms give
\[
 \wp(a_s)=\wp(\delta_s)-(s(\gamma)-\gamma)=0.
\]
The roots of $X^p-X$ in $K$ are exactly $\mathbf F_p$:
the $p$ prime-field elements are distinct roots of this
degree-$p$ polynomial, so there can be no others.
Thus each $a_s$ lies in $\mathbf F_p$ and is fixed by $G$.
In particular $a_1=0$.
Subtracting the identity
\[
 (s(\beta)-\beta)+s(t(\beta)-\beta)
       -(st(\beta)-\beta)=0
\]
from the equation for $\delta$ gives
$a_s+a_t-a_{st}=c(s,t)$.

Define $v_s=\iota(-a_s)u_s$.
The central-coordinate multiplication formula now gives
\[
 v_sv_t
   =\iota(-a_s-a_t+c(s,t))u_{st}
   =\iota(-a_{st})u_{st}=v_{st}.
\]
Also $v_1=1$ and $q(v_s)=s$, so $v$ is a group-homomorphic
section of $q$. Its image $H$ is a subgroup of $E$.
We have $NH=E$, since replacing $u_s$ by $v_s$ changes it
by an element of $N$. The Frattini lemma forces $H=E$.
On the other hand, $H\cap N=\{1\}$: if $v_s\in N$, then
$s=q(v_s)=1$ and $v_s=v_1=1$.
This contradicts $H=E$ and $|N|=p>1$.
\end{proof}

\begin{lemma}[Semilinear automorphisms of the Artin--Schreier quotient]
\label{lem:witt-semilinear-quotient-lifts}
\uses{thm:witt-additive-equations}
\source{witt1936:page-0004,camina1997:page-0004}
\lean{Camina1997.Witt.liftEquiv, Camina1997.Witt.wittLift_algebraMap, Camina1997.Witt.wittLift_root, Camina1997.Witt.wittLift_mul, Camina1997.Witt.wittLift_one, Camina1997.Witt.wittLift_injective, Camina1997.Witt.wittLiftAlgEquiv}
\leanok
Suppose a group $G$ acts faithfully on a field $K$ of characteristic
$p$. Let $c:G^2\to\mathbf F_p$, $\delta:G\to K$, and $\gamma\in K$
satisfy
\[
 \delta_1=0,\qquad
 \delta_s+s(\delta_t)-\delta_{st}=c(s,t),\qquad
 s(\gamma)-\gamma=\delta_s^p-\delta_s.
\]
Put $A=K[Y]/(Y^p-Y-\gamma)$ and let $y$ denote the class of $Y$.
For each $s\in G$ and $a\in\mathbf F_p$ there is a ring automorphism
$\varphi_{s,a}$ of $A$ satisfying
\[
 \varphi_{s,a}(x)=s(x)\quad(x\in K),\qquad
 \varphi_{s,a}(y)=y+\delta_s+a.
\]
The map $(s,a)\mapsto\varphi_{s,a}$ is injective, and
\[
 \varphi_{s,a}\varphi_{t,b}=\varphi_{st,a+b+c(s,t)},\qquad
 \varphi_{1,0}=1.
\]
If the action on $K$ fixes a coefficient algebra $F$, these are
$F$-algebra automorphisms. No irreducibility of the modulus is needed here.
\end{lemma}

\begin{proof}
More generally, if $d^p-d=s(\gamma)-\gamma$, apply $s$ to
coefficients and send $Y$ to $y+d$. The defining polynomial maps to
\[
 (y+d)^p-(y+d)-s(\gamma)
 =\gamma+(d^p-d)-s(\gamma)=0,
\]
so this map descends to $A$. Its inverse applies $s^{-1}$ to
coefficients and sends $y$ to $y-s^{-1}(d)$; the inverse correction
satisfies the same relation for $s^{-1}$. The composites fix the
coefficients and $y$, hence every residue class.
Take $d=\delta_s+a$, using $a^p=a$.
The formulas on $K$ and $y$ give the asserted multiplication rule
by the equation for $\delta$, and give the identity when $s=1,a=0$.
The coefficient map $K\to A$ is injective because the monic modulus
has positive degree. Equality of two such automorphisms therefore
implies equality of their actions on $K$, whence $s=t$ by
faithfulness. Their values on $y$ then give $a=b$.
Finally the coefficient formula shows that each automorphism fixes
$F$ whenever the original action does.
\end{proof}

\begin{lemma}[A faithful central-extension action on a quotient ring]
\label{lem:witt-central-quotient-action}
\uses{lem:central-extension-factor-set,thm:witt-additive-equations,lem:witt-semilinear-quotient-lifts}
\source{witt1936:page-0004,camina1997:page-0004}
\lean{Camina1997.Witt.exists_centralLift, Camina1997.Witt.centralLift, Camina1997.Witt.centralLift_injective, Camina1997.Witt.centralLift_algebraMap, Camina1997.Witt.centralLift_root, Camina1997.Witt.centralLiftAlgHom, Camina1997.Witt.centralLiftAlgHom_injective}
\leanok
More generally, let a finite group $G$ act faithfully on a field $K$
of characteristic $p$, and let
\[
 1\longrightarrow (\mathbf F_p,+)\xrightarrow{\iota}E
   \xrightarrow{q}G\longrightarrow1
\]
be a central extension. There are $\gamma\in K$ and an injective
homomorphism $\rho:E\to\operatorname{Aut}(A)$, where
$A=K[Y]/(Y^p-Y-\gamma)$, such that
\[
 \rho(e)(x)=q(e)(x)\qquad(x\in K).
\]
If the action on $K$ fixes a coefficient algebra $F$, the same
construction takes values in $\operatorname{Aut}_F(A)$.
This assertion concerns the quotient ring; it does not assert
that the quotient is a field.
\end{lemma}

\begin{proof}
Choose a normalized section $u$ of $q$ and its factor set $c$.
The trace-one formulas construct $\delta,\gamma$ satisfying
Witt's equations for $c$, with $\delta_1=0$.
Assign to $e=\iota(a)u_s$ the automorphism $\varphi_{s,a}$ of
the preceding lemma. Unique central coordinates make this a
well-defined function. The central-coordinate multiplication
formula agrees exactly with
$\varphi_{s,a}\varphi_{t,b}=\varphi_{st,a+b+c(s,t)}$,
and the identity corresponds to $(1,0)$, so it is a homomorphism.
Injectivity follows from injectivity of the coordinate equivalence
and of $(s,a)\mapsto\varphi_{s,a}$. The coefficient formula gives
both the restriction identity and the assertion about $F$.
\end{proof}

\begin{lemma}[The lifted action on the Artin--Schreier field]
\label{lem:witt-lifted-automorphisms}
\uses{lem:central-extension-factor-set,lem:witt-frattini-parameter-not-image,lem:artin-schreier-irreducible,lem:witt-semilinear-quotient-lifts}
\source{witt1936:page-0004,camina1997:page-0004}
\notready
For $\delta$ and $\gamma$ satisfying the equations of
Lemma~\ref{lem:witt-frattini-parameter-not-image}, set
\[
 L=K[Y]/(Y^p-Y-\gamma)
\]
and let $\theta$ be the class of $Y$.
Then $L$ is a field with $[L:K]=p$, and there is an injective
homomorphism $\rho:E\to\operatorname{Aut}_F(L)$ given by
\[
 \rho(\iota(a)u_s)(x)=s(x)\quad(x\in K),\qquad
 \rho(\iota(a)u_s)(\theta)=\theta+a+\delta_s.
\]
\end{lemma}

\begin{proof}
Lemma~\ref{lem:witt-frattini-parameter-not-image} gives
$\gamma\notin\wp(K)$.
The Artin--Schreier irreducibility lemma therefore applies to
$Y^p-Y-\gamma$. Its quotient is a field: a nonzero residue
class, represented by a polynomial of degree less than $p$,
is relatively prime to the irreducible modulus, so a Bezout
identity supplies its inverse. Monic division also gives the
basis $1,\theta,\ldots,\theta^{p-1}$ over $K$ and degree $p$.

For $a\in\mathbf F_p$ and $s\in G$, apply $s$ to polynomial
coefficients and evaluate the variable at $\theta+a+\delta_s$.
This defines a unital $F$-algebra map $K[Y]\to L$.
The image of its modulus is zero, because
\[
 (\theta+a+\delta_s)^p-(\theta+a+\delta_s)
   =\gamma+\wp(\delta_s)=s(\gamma).
\]
Thus it factors through an $F$-algebra map $\varphi_{a,s}:L\to L$
with the displayed formulas. These formulas determine the map,
since every element is a polynomial in $\theta$ with
coefficients in $K$.

On $K$, the composite $\varphi_{a,s}\varphi_{b,t}$ acts by $st$.
On the root it acts by
\[
 \theta+a+b+\delta_s+s(\delta_t)
   =\theta+a+b+c(s,t)+\delta_{st}.
\]
Consequently
$\varphi_{a,s}\varphi_{b,t}
 =\varphi_{a+b+c(s,t),st}$.
Also $\varphi_{0,1}$ is the identity because $\delta_1=0$.
The coordinate multiplication in $E$ shows that assigning
$\varphi_{a,s}$ to $\iota(a)u_s$ respects multiplication and
identity. In particular the element assigned to $e^{-1}$
is a two-sided inverse of the element assigned to $e$.
These maps are therefore automorphisms and define $\rho$.

If $\rho(\iota(a)u_s)$ is the identity, restriction to $K$
gives $s=1$. Its action on $\theta$ then gives
$\theta+a=\theta$, so $a=0$. Thus $\rho$ has trivial kernel
and is injective.
\end{proof}

\begin{lemma}[The lifted action is the full Galois group]
\label{lem:witt-lift-full-galois-group}
\uses{lem:witt-lifted-automorphisms}
\source{witt1936:page-0004,camina1997:page-0004}
\notready
The field $L$ above is a finite Galois extension of $F$.
The map $\rho$ is an isomorphism
$E\simeq\operatorname{Gal}(L/F)$, and restriction to $K$
under this isomorphism is exactly $q$.
\end{lemma}

\begin{proof}
The extension $L/F$ is finite, since
$[L:F]=[L:K][K:F]=p[K:F]$.
It is separable: $K/F$ is separable, and the minimal polynomial
$Y^p-Y-\gamma$ of $\theta$ over $K$ has derivative $-1$;
separability is transitive in a finite tower.

Choose an algebraic closure containing $L$ and let
$j:L\to\overline F$ be any $F$-embedding.
Normality of $K/F$ implies $j(K)=K$, so its restriction is
some $s\in G$. The element $j(\theta)$ is a root of
$Y^p-Y-s(\gamma)$. That polynomial has the $p$ distinct roots
\[
 \theta+\delta_s+a\qquad(a\in\mathbf F_p),
\]
because $\wp(\theta+\delta_s)=s(\gamma)$; these are all its
roots by its degree. Hence $j(\theta)\in L$, and the
polynomial basis shows $j(L)\subseteq L$.
This proves normality of $L/F$. Explicitly, for $z\in L$
and any root $z'$ of its minimal polynomial over $F$, the
embedding $F(z)\to\overline F$ sending $z$ to $z'$ extends
to $L$: adjoin a finite list of algebraic generators and at
each step send the next generator to a root of its transported
minimal polynomial in the algebraic closure. The resulting
$F$-embedding sends $L$ into $L$ by the argument just given,
so $z'\in L$. Thus every such minimal polynomial splits in
$L$, the defining normality condition.

Finally let $\tau\in\operatorname{Gal}(L/F)$. Its restriction
to $K$ is some $s\in G$, and the same root calculation gives
$\tau(\theta)=\theta+\delta_s+a$ for a unique
$a\in\mathbf F_p$.
Therefore $\tau$ and $\rho(\iota(a)u_s)$ agree on $K$
and $\theta$, hence on all of $L$. This proves surjectivity;
injectivity was proved in the preceding lemma.
The displayed action on $K$ says
$\rho(e)|_K=q(e)$ for every $e\in E$, as asserted.
\end{proof}

\begin{theorem}[Proper lifting across a central Frattini kernel of order $p$]
\label{thm:witt-central-frattini-lift}
\uses{lem:central-extension-factor-set,thm:witt-additive-equations,lem:witt-frattini-parameter-not-image,lem:witt-lifted-automorphisms,lem:witt-lift-full-galois-group}
\source{witt1936:page-0003,witt1936:page-0004,camina1997:page-0003,camina1997:page-0004}
\notready
Let $F$ have characteristic $p$, let $K/F$ be finite Galois,
and let $E$ be a finite $p$-group. Suppose
\[
 q:E\longrightarrow\operatorname{Gal}(K/F)
\]
is surjective and $N=\ker q$ has order $p$ with
$N\le Z(E)\cap\Phi(E)$.
There is a field extension $L/K$ of degree $p$ such that
$L/F$ is Galois and an isomorphism
\[
 \rho:E\longrightarrow\operatorname{Gal}(L/F)
 \quad\text{satisfying}\quad\rho(e)|_K=q(e).
\]
\end{theorem}

\begin{proof}
Choose a nonidentity element of $N$. Its order divides $p$
and is not one, so it has order $p$ and generates $N$.
Sending $a\in\mathbf F_p$ to the corresponding power of this
generator gives an isomorphism $\iota:(\mathbf F_p,+)\to N$.
Choose a representative $u_s$ of each fiber of $q$, taking
$u_1=1$. The factor-set lemma supplies a normalized
$\mathbf F_p$-valued two-cocycle $c$.
The theorem on Witt's additive equations constructs
$\delta_s,\gamma\in K$ satisfying the three required
identities. The Frattini parameter lemma proves
$\gamma\notin\wp(K)$.
Form the actual quotient field
$L=K[Y]/(Y^p-Y-\gamma)$ and the lifted automorphisms as above.
The quotient has degree $p$ over $K$, and the full-Galois-group
lemma gives the required isomorphism and its restriction
identity.
\end{proof}
